% ---
% id: reductions/domination/thm:noise-channel
% kind: thm
% title: "The noise channel of a domination"
% summary: >-
%   If G dominates F sector by sector, G squashes exactly to F with added flag
%   sectors, or to F itself where F has flag sectors, through a noise channel that keeps
%   each sector with probability one minus its weight. The adjoint returns each
%   sector projector scaled by that probability.
% prominence: main
% status: accepted
% uses: [def:domination-link, foundations/def:domination, foundations/lem:domination, reductions/def:link, reductions/def:squashing, reductions/def:generalized-instrument, reductions/lem:master, reductions/prop:common, reductions/lem:top]
% mentions: [reductions/rem:dominating]
% assumes: []
% refs:
%   - "nahar2025imperfect, Thm. 3, Eqs. (20)-(24)"
%   - "nahar2025imperfect, Technical aside 1"
% provenance: newly-derived
% lean:
% notes: "The published theorem states case (a), flags in a flag space added for the purpose, on a block decomposition of the physical space, with q_m in [0,1]; here it is derived from the branching lemma and stated on a finite target system, with the shared-flag case (b) added and the weights restricted to [0,1), where the residual family of foundations/def:domination is defined. No requirement uses the theorem; it is kept for a declared device weight (A-DOM, spec/future-scope.md section 13)."
% ---

The channel is the one of \cite[Thm.~3, Eqs.~(20)--(24)]{nahar2025imperfect},
derived here from the branching calculus of this contract: a \tref{generalized
instrument} decides, sector by sector, whether to keep the input
(probability $1-\vartheta_s$) or to flag it (probability $\vartheta_s$); the keep
branch passes the input on unchanged, and the flag branch measures the
residual family and prepares the flag of its outcome. The domination
$G=(1-\vartheta)F+\vartheta R$ sector by sector is exactly the statement that the
two branches together reproduce $G$. Remark~\ref{rem:dominating} reads the
same split as a dominated detector. Part~(iii), in which the flags land in
sectors the target already has, is the use of an existing flag space that
\cite[Technical aside~1]{nahar2025imperfect} describes.

\begin{theorem}\label{thm:noise-channel}
In the setting of Definition~\ref{def:domination-link}, let $G$ dominate $F$
at the \tref{domination weights} $\vartheta$ and let $R$ be the residual family.
\begin{enumerate}[label=(\roman*),leftmargin=3.0em]
\item $R=\{R_v\}_{v\in\Vset}$ is a POVM on $\Qsys$, block diagonal for
  $\{\Pi_s\}$, and
  $G^{(s)}_v=(1-\vartheta_s)F^{(s)}_v+\vartheta_sR^{(s)}_v$ for every $v$ and $s$.
\item \emph{(added flag sectors)} $\Phi^{+}_{\vartheta}$ of \eqref{eq:noise-fresh} is a
  channel $\Tcl{\Qsys}\to\Tcl{\Qsys^{+}}$, and
  \begin{equation}\label{eq:noise-fresh-relation}
    \big(G;\,(1-\vartheta_s)\Pi_s\,(s\in\SecSet),\,\textstyle\sum_s\vartheta_s\Pi_s\big)
    \ \dsq_{\Phi^{+}_{\vartheta}}\
    \big(F^{+};\,\Pi_s\dsum0\,(s\in\SecSet),\,0\dsum\idon{\Flag^{+}}\big).
  \end{equation}
\item \emph{(shared flag sectors)} If $F$ has \tref{indicator sectors} $(f)$ with
  $\vartheta_{(f)}=0$ for every $f$, then
  $\Phi_{\vartheta}$ of \eqref{eq:noise-shared} is a channel
  $\Tcl{\Qsys}\to\Tcl{\Qsys}$ with $G\dsq_{\Phi_{\vartheta}}F$, and its adjoint
  takes the sector projectors and the flag projector to
  \begin{equation}\label{eq:noise-shared-pullbacks}
  \begin{gathered}
    \Phi_{\vartheta}^{\dagger}(\Pi_s)=(1-\vartheta_s)\Pi_s\quad
    (s\text{ not an indicator sector}),\\
    \Phi_{\vartheta}^{\dagger}(\Pi_{(f)})=\Pi_{(f)}+\sum_{s}\vartheta_s\,
      \Pi_sR_f\Pi_s,\qquad
    \Phi_{\vartheta}^{\dagger}(\Pflag)=\Pflag+\sum_{s}\vartheta_s\Pi_s .
  \end{gathered}
  \end{equation}
\item \emph{(output)} The outputs of both channels are block diagonal for the
  sectors of their target, so either channel is unchanged when the sector
  pinching of its target is applied after it.
\item \emph{(weight zero)} If $\vartheta_s=0$ for every $s$, then $F=G$, and
  $\Phi_{\vartheta}$, where defined, is the sector pinching
  $\sigma\mapsto\sum_s\Pi_s\sigma\Pi_s$.
\end{enumerate}
In particular the domination link of Definition~\ref{def:domination-link}, in
either case, is a receiver-stage link of Definition~\ref{def:link}: its
channel witness relates its source family and carried operators to its
target and carried operators exactly, as \eqref{eq:link} requires with
$\textsf{status}=\textsf{exact}$.
\end{theorem}

\begin{proof}
\emph{Idea.} The two channels are the composite channels of
Lemma~\ref{lem:master} and Proposition~\ref{prop:common} for one generalized
instrument with a keep branch and a flag branch.

\emph{(i)} On every sector $s$, $G^{(s)}$ and $F^{(s)}$ are Hermitian, positive
semidefinite and complete at $1$, $G$ and $F$ being POVMs block diagonal for
$\{\Pi_s\}$. \xref{foundations/lem:domination}\,(vi) gives that $R^{(s)}$ has
positive semidefinite elements, is complete at $1$ and satisfies the stated
split; so $R_v=\bigoplus_sR^{(s)}_v\succeq0$ and
$\sum_vR_v=\bigoplus_s\idon{d_s}=\idon{\Qsys}$.

\emph{The instrument.} Put
$\Ins_{\mathrm{k}}(\sigma)\coloneqq\sum_s(1-\vartheta_s)\Pi_s\sigma\Pi_s$ and
$\Ins_{\mathrm{f}}(\sigma)\coloneqq\sum_s\vartheta_s\Pi_s\sigma\Pi_s$, both maps
$\Tcl{\Qsys}\to\Tcl{\Qsys}$. Each is completely positive, a sum of maps
$\sigma\mapsto c\,\Pi_s\sigma\Pi_s$ with $c\ge0$, and
$\tr\Ins_{\mathrm{k}}(\sigma)+\tr\Ins_{\mathrm{f}}(\sigma)
=\sum_s\tr[\Pi_s\sigma\Pi_s]=\tr\sigma$, since $\sum_s\Pi_s=\idon{\Qsys}$.
So $\{\Ins_{\mathrm{k}},\Ins_{\mathrm{f}}\}$ is a generalized instrument
(Definition~\ref{def:generalized-instrument}). Its adjoints are
$\Ins_{\mathrm{k}}^{\dagger}(X)=\sum_s(1-\vartheta_s)\Pi_sX\Pi_s$ and
$\Ins_{\mathrm{f}}^{\dagger}(X)=\sum_s\vartheta_s\Pi_sX\Pi_s$. In the keep branch
take the POVM $F$ on $\Qsys$ with the identity channel, so
$F\dsq_{\idmap}F$; in the flag branch take the POVM $R$ on $\Qsys$ with the
measure-and-prepare channel
$\Lmp(\tau)=\sum_v\tr[R_v\tau]\dyad{v}_{\Flag^{+}}$, so $R\dsq_{\Lmp}\Ppvm$ by
Lemma~\ref{lem:top}. By Lemma~\ref{lem:master}\,(i) the branch elements are
$M_{v,\mathrm{k}}=\Ins_{\mathrm{k}}^{\dagger}(F_v)=\sum_s(1-\vartheta_s)\Pi_sF_v\Pi_s$
and $M_{v,\mathrm{f}}=\Ins_{\mathrm{f}}^{\dagger}(R_v)=\sum_s\vartheta_s\Pi_sR_v\Pi_s$,
and since $F_v$ and $R_v$ are block diagonal,
$M_{v,\mathrm{k}}+M_{v,\mathrm{f}}
=\bigoplus_s\big((1-\vartheta_s)F^{(s)}_v+\vartheta_sR^{(s)}_v\big)=G_v$ by (i).

\emph{(ii)} Lemma~\ref{lem:master}\,(iii) gives
$\{M_{v,\mathrm{k}}+M_{v,\mathrm{f}}\}_v\dsq_{\Lambda}\{F_v\dsum\dyad{v}\}_v$
with $\Lambda(\sigma)=\Ins_{\mathrm{k}}(\sigma)\dsum\Lmp\big(\Ins_{\mathrm{f}}(\sigma)\big)$,
a channel by Lemma~\ref{lem:master}\,(ii). The left side is $G$, the right side
is $F^{+}$, and
$\Lmp(\Ins_{\mathrm{f}}(\sigma))=\sum_v\sum_s\vartheta_s\tr[R_v\Pi_s\sigma\Pi_s]\dyad{v}$,
so $\Lambda=\Phi^{+}_{\vartheta}$ of \eqref{eq:noise-fresh}. For the carried
operators, \eqref{eq:masteradjoint} gives
$(\Phi^{+}_{\vartheta})^{\dagger}(X\dsum Y)=\Ins_{\mathrm{k}}^{\dagger}(X)
+\Ins_{\mathrm{f}}^{\dagger}\big(\Lmp^{\dagger}(Y)\big)$, with
$\Lmp^{\dagger}(Y)=\sum_v\bra{v}Y\ket{v}R_v$. At $X=\Pi_s$, $Y=0$ this is
$(1-\vartheta_s)\Pi_s$; at $X=0$, $Y=\idon{\Flag^{+}}$ it is
$\Ins_{\mathrm{f}}^{\dagger}(\sum_vR_v)=\Ins_{\mathrm{f}}^{\dagger}(\idon{\Qsys})
=\sum_s\vartheta_s\Pi_s$, which is Lemma~\ref{lem:master}\,(iv) for the flag
branch. That is \eqref{eq:noise-fresh-relation} in the notation of
Definition~\ref{def:squashing}.

\emph{(iii)} Let $W\colon\mathbb{C}^{|\Vset|}\to\Hs{\Qsys}$ be the isometry with
$W\ket{f}$ the unit vector spanning the indicator sector $(f)$, and
$\Psi(\tau)\coloneqq W\tau W^{\dagger}$, a channel
$\Tcl{\Flag^{+}}\to\Tcl{\Qsys}$ with adjoint $\Psi^{\dagger}(X)=W^{\dagger}XW$.
Because every $F_v$ is block diagonal and equals $\delta_{vf}$ on the sector
$(f)$, $\Psi^{\dagger}(F_v)=\sum_f\delta_{vf}\dyad{f}=\dyad{v}$, so
$\Ppvm\dsq_{\Psi}F$. With $F\dsq_{\idmap}F$ in the keep branch, $F$ is a
common target of the two branch targets $F$ and $\Ppvm$, and
Proposition~\ref{prop:common} gives
$\{M_{v,\mathrm{k}}+M_{v,\mathrm{f}}\}_v\dsq_{\Lambda'}F$ with
$\Lambda'=\idmap\circ\idmap\circ\Ins_{\mathrm{k}}+\Psi\circ\Lmp\circ\Ins_{\mathrm{f}}$,
a channel by that proposition. The left side is $G$, and
$\Psi(\Lmp(\Ins_{\mathrm{f}}(\sigma)))=\sum_f\sum_s\vartheta_s\tr[R_f\Pi_s\sigma\Pi_s]\,W\dyad{f}W^{\dagger}$
with $W\dyad{f}W^{\dagger}=\Pi_{(f)}$, so $\Lambda'=\Phi_{\vartheta}$ of
\eqref{eq:noise-shared}. Its adjoint is
$\Phi_{\vartheta}^{\dagger}(X)=\Ins_{\mathrm{k}}^{\dagger}(X)
+\Ins_{\mathrm{f}}^{\dagger}\big(\sum_f\bra{f}W^{\dagger}XW\ket{f}R_f\big)$,
by linearity and the adjoints of the three maps. For a sector $s$ that is not
an indicator sector, $W^{\dagger}\Pi_sW=0$, so
$\Phi_{\vartheta}^{\dagger}(\Pi_s)=(1-\vartheta_s)\Pi_s$. For an indicator sector,
$\vartheta_{(f)}=0$ gives $\Ins_{\mathrm{k}}^{\dagger}(\Pi_{(f)})=\Pi_{(f)}$, and
$W^{\dagger}\Pi_{(f)}W=\dyad{f}$ gives the second term
$\Ins_{\mathrm{f}}^{\dagger}(R_f)=\sum_s\vartheta_s\Pi_sR_f\Pi_s$. Summing over
$f$, with $\sum_fR_f=\idon{\Qsys}$ by (i),
$\Phi_{\vartheta}^{\dagger}(\Pflag)=\Pflag+\sum_s\vartheta_s\Pi_s$.

\emph{(iv)} Every term of \eqref{eq:noise-fresh} and \eqref{eq:noise-shared}
is $\Pi_s\sigma\Pi_s$ for a sector $s$, a multiple of $\dyad{v}_{\Flag^{+}}$, the
projector onto the added sector $(v)^{+}$ of $\Qsys^{+}$, or a multiple
of $\Pi_{(f)}$; each is block diagonal for the sectors of the target, and the
sector pinching fixes it.

\emph{(v)} If every $\vartheta_s=0$, \eqref{eq:dom-link} reads
$G^{(s)}_v-F^{(s)}_v\succeq0$ for all $v$ and $s$, and both families are
complete at $1$ on every sector, so $F^{(s)}=G^{(s)}$ by
\xref{foundations/lem:domination}\,(iii). The second sum of
\eqref{eq:noise-shared} vanishes and the first is the sector pinching.

\emph{The closing sentence.} The carried operators of the link in case (a)
are those of \eqref{eq:noise-fresh-relation}, and in case (b) the
projectors of the sectors that are not indicator sectors and $\Pflag$, whose
pull-backs are \eqref{eq:noise-shared-pullbacks}; so (ii) and (iii) are the
relation the channel-witness slot of \eqref{eq:link} asks for, with no error
term.
\end{proof}
